% INTEGRATED into 02_paper1.tex on 17 Aug 2026; kept for the record, not \input anywhere.
% ============================================================
% REWRITE DRAFT — Chapter 2, sections 1 and 2, in the economics-paper
% register (thesis/WRITING_STYLE.md, 16 Aug 2026). For Damian's check of
% direction before sections 3-5 are written. Numbers are macros where a
% macro exists; the rest convert at integration.
% ============================================================

\section{Introduction}
\label{sec:p1_intro}

School culture is widely held to matter for what pupils learn, but it is
rarely measured. This paper constructs measures of two dimensions of the
culture of English secondary schools, warmth and strictness, for the full
population of around 3,330 state secondary schools, and asks how far each
dimension can be recovered from data that exist for every school.

Warmth is the relational quality of a school: whether pupils are known,
cared for and supported by the adults around them. Strictness is its
behavioural structure: whether expectations of conduct are clear, enforced
and met. The two are the axes of the parenting typology of
\citet{baumrind1971} and \citet{maccoby1983}, in which the combination of
high responsiveness with high demandingness, the authoritative style,
produces the best child outcomes; the same two-dimensional structure has been
applied to schools by \citet{pellerin2005} and \citet{gregory2010}, who find
that pupils in schools combining support with structure report less
victimisation and more engagement. Warmth and strictness are treated here as
distinct dimensions rather than the ends of one axis, because a school can be
warm without being orderly and orderly without being warm, and the empirical
question of whether the two co-vary is left to the data. The difficulty is
that neither is recorded in any administrative dataset. What is
recorded, for every school, is text: an inspection report written by an
external observer, a behaviour policy written by the school, and a website
written by the school for parents. What is recorded for a smaller set of
schools, because I collected it, is direct observation and a structured
interview with the headteacher.

The design follows from that asymmetry. For \NTierOne{} schools I have a
research visit, comprising structured lesson and out-of-lesson observations,
and a headteacher interview; from these I construct a gold-standard measure of
each dimension. Because the interview and the visit disagree substantially
about the same school, the gold standard is two measures, not one: an
\emph{enacted} score, from what an observer recorded, and an \emph{espoused}
score, from what the headteacher reported. For a further 201 schools I have
the interview alone, and for the remaining national population I have only
the public text. For those schools I estimate the gold standard from what is
observable, using prediction models fitted on the schools where both the
predictors and the gold standard exist, and I report how well the predictions
reproduce the gold standard out of sample.

Three findings follow. First, the two dimensions are recoverable to very
different degrees. Strictness scored from inspection-report text is
correlated with observed strictness at $r = \CorrOfstedStrictness$, and a
text-only prediction model retains a leave-one-out correlation of
$\RidgeBStrictness$ with the gold standard. Warmth is not recoverable from
any public source: no text-derived warmth score is correlated with observed
warmth, and the text-only model has no out-of-sample skill (leave-one-out
$R^2$ of $-0.02$). Second, the enacted and espoused measures diverge in a
systematic and dimension-specific way. Headteachers rate their schools nearly
two points warmer on a ten-point scale than observers record and about one
point stricter; the interview and
the visit are correlated at $\GoldWarmthSplit$ for warmth and
$\GoldStrictnessSplit$ for strictness. Every public source, including a
survey of parents, tracks the espoused warmth measure and not the enacted one.
Third, the language-model instruments used to score the text carry a
measurement result of their own: an instrument's agreement with its reference
labels and its correlation with the observed measure are unrelated, and the
best-agreeing version of an instrument was usually the worse measure.

The paper contributes a measurement framework and a set of school-level
scores, a national-scale demonstration of what school culture is and is not
legible from public text, and a set of findings about language-model
measurement that apply beyond this setting. The scores are the explanatory
variables of the following chapter.

The paper sits between three literatures. The school-climate literature
measures culture through pupil, teacher and parent surveys
\citep{hoy2001,cohen2009,thapa2013}; these capture perceptions, are open to
social desirability, and are costly at national scale, and administrative
proxies such as exclusion rates \citep{machin2020} conflate culture with
intake. Structured classroom observation \citep{hamre2007} measures practice
directly but has not, to my knowledge, been combined with a leadership
interview and public text to separate what a school does from what it says.
And the text-as-data literature in economics \citep{gentzkow2019} supplies
the tools for reading the public record at scale; this paper applies large
language models rather than dictionary or topic methods, and reports what
that choice costs and buys in measurement terms. Relative to all three, the
contribution is a measure of enacted culture with a national text-based
extension whose validity is tested against direct observation rather than
assumed.

The paper proceeds as follows. Section~2 describes the data. Section~3
constructs the gold standard, the text-based instruments and the prediction
models, and reports how well the predictions reproduce the gold standard.
Section~4 reports what the measures show. Section~5 concludes.

% ============================================================
\section{Data}
\label{sec:p1_data}

The paper uses three kinds of data on English state secondary schools: a
research visit and a headteacher interview conducted for this study, three
corpora of public text, and the Department for Education's school-level
administrative panel. \Cref{tab:tier_summary} sets out coverage.

\textbf{Visits and interviews.} Between November 2023 and April 2025 I
interviewed the headteacher, or where unavailable the senior leader
responsible for behaviour and culture, at 304 schools, and visited
\NTierOne{} of them. Schools were recruited by direct approach and by
referral from participating schools; no quota was applied. Three-quarters of
interviews and half of visits took place in the 2024/25 academic year.

A visit comprised structured observation of lessons and of unstructured time,
each by two researchers rating independently on a fixed schedule. The lesson
schedule has seven items bearing on warmth and strictness (for example the
frequency with which the teacher uses pupils' names; the level of misbehaviour
tolerated before a first sanction) and nine on teaching practice; the
out-of-lesson schedule has eight items covering transitions, breaks, corridors
and the dining hall. All items are rated 1 to 5 and averaged across
observations. The interview followed a fixed schedule of eleven questions and
nine statements on which the headteacher rated their own school on a five-point
agreement scale (for example ``Sanction systems are applied consistently across
the school''; ``Pupils in this school feel that their teachers care about
them''), and recorded whether eight named behaviour systems were in place. The
full schedules are reproduced in Appendix~2A. Inter-rater agreement on the
observation items is reported with the measures in Section~3.

\textbf{Public text.} Three corpora cover the national population. Ofsted
inspection reports were obtained for 3,288 of the 3,312 open schools
(99.3 per cent); the remainder opened after 2023 and had not been inspected.
Behaviour policy documents were obtained for 3,365 schools (98.2 per cent),
one document per school, selected by a rule that prefers the current
whole-school behaviour policy over annexes and predecessor versions. School
websites were crawled once for 3,053 schools, taking the home, values and
welcome pages; 3,370 crawl files exist including partial ones. Where a crawl
returned no prose, the school is flagged rather than dropped.

\textbf{Administrative data.} All schools are matched on their unique
reference number to a school-year panel built from the Department for
Education's published tables for 2015/16 to 2024/25, carrying pupil
composition, prior attainment, Progress~8 and its components, workforce and
finance. This paper uses the panel only to describe the samples; the
following chapter uses it for outcomes and controls.

\textbf{Representativeness.} \Cref{tab:2A:representativeness} compares the
visited and interviewed schools with the national population. Both samples
resemble the population on school type, phase, admissions, size, deprivation
and urban location, and both over-represent London (45 per cent of visited
schools against 15 per cent nationally). Both also lean toward stronger
schools: 25 per cent of visited schools held an Outstanding inspection grade
as at August 2024 against 15 per cent nationally, and their mean Progress~8
in 2018/19 was $+0.17$ against $+0.04$. Estimates from the visited sample
therefore apply most directly to a population somewhat stronger than the
national average, and the restricted range of both culture and outcomes in
that sample works against finding associations rather than for it.

\Cref{tab:score_dist} reports descriptive statistics for the measures
constructed in Section~3, by sample.

\begin{table}[htbp]
\centering
\small
\begin{tabular}{lcccccc}
\toprule
Group & $N$ & FSM\% & EAL\% & Median size & Outstanding\% & Good\% \\
\midrule
Visited and interviewed & 103 & 29.2 & 27.8 & 1{,}118 & 25 & 65 \\
Interviewed only & 200 & 27.2 & 17.9 & 1{,}052 & 15 & 75 \\
Rest of the national population & 3{,}028 & 30.1 & 18.2 & 1{,}050 & 14 & 70 \\
\midrule
All schools & 3{,}332 & 29.9 & 18.4 & 1{,}051 & 15 & 70 \\
\bottomrule
\end{tabular}
\caption{Distribution of schools across the three groups and selected observable characteristics (2023--24 school year). The groups are mutually exclusive: the first received both a headteacher interview and a researcher visit, the second an interview only, and the third is the rest of the national population. FSM\%: mean proportion eligible for free school meals.
EAL\%: mean proportion with English as an additional language. Size: median
number of pupils on roll. Outstanding\%/Good\%: percentage rated Outstanding or
Good at most recent inspection, among schools with a grade.}
\label{tab:tier_summary}
\smallskip
\begin{minipage}{0.92\linewidth}
\footnotesize\textit{Notes:} The groups sum to one fewer than the total: one interviewed school whose head left the statement battery blank belongs to the total but to no group.
\end{minipage}
\end{table}

